
\chapter{Categories, Functors and Natural Transformations}

\section{Categories}

It is quite easy to make examples motivating the definition of
categories and the evolution that follows through these pages.

\begin{example}[Set Theory]\label{example:Set}
  Here, we have {\em sets} and {\em functions}. Whereas the concepts of
  set ad membership are primitive, functions are formalised as follows:
  for \(A\) and \(B\) sets, a function from \(A\) to \(B\) is any
  \(f \subseteq A \times B\) such that for every \(x \in A\) there exists one and only
  one \(y \in B\) such that \((x, y ) \in f\). We write
  \[
    f : A \to B \ \text{or} \ A \functo f B
  \]
  to say \q{\(f\) is a function from \(A\) to \(B\)} and, for
  \(x \in A\), we write \(f(x)\) the element of \(B\) bound to \(x\) by
  \(f\). Consecutive functions can be combined in a quite natural way:
  for \(A\), \(B\) and \(C\) sets and functions
  \[
    A \functo{f} B \functo {g} C ,
  \]
  the {\em composite} of \(g\) and \(f\) is the function
  \[
    g \circ f : A \to C \,,\ g \circ f(x) := g\big(f(x)\big) .
  \]
  Informally speaking: \(f\) takes one input and gives one output; it is
  then passed to \(g\), which then provides one result. Such operation
  is called {\em composition} and has some nice basic properties
  \begin{tcbenum}
  \item Every set \(A\) has associated an {\em identity}
    \[
      \id_A : A \to A \,,\ \id_A (x) := x
    \]
    is such that for every set \(B\) and function \(g : B \to A\) we have
    \[
      \id_A \circ g = g
    \]
    and for every set \(C\) and function \(h : A \to C\) we have
    \[
      h \circ \id_A = h .
    \]
  \item \(\circ\) is {\em associative}, that is for \(A\), \(B\), \(C\) and
    \(D\) sets and
    \[
      A \functo{f} B \functo{g} C \functo{h} D
    \]
    functions, we have the identity
    \[
      (h \circ g) \circ f = h \circ (g \circ f) .
    \]
  \end{tcbenum}
\end{example}

\begin{example}[Topology]\label{example:Top}
  A {\em topological space} is a set where some of its subsets have the
  status of \q{open} sets. Being sets at their core, we have functions
  between topological spaces, but some of them are more interesting than
  others. Namely, {\em continuous functions} are functions that care
  about the label of open: for if \(X\) and \(Y\) are topological
  spaces, a function \(f : X \to Y\) is said {\em continuous} whenever for
  every open set \(U\) of \(Y\) the set \(\inv f U\) is an open set of
  \(X\). Being function, consecutive continuous functions can be
  composed: is the resulting function continuous as well? Yes: if \(X\),
  \(Y\) and \(Z\) are topological spaces and \(f\) and \(g\) continuous,
  for if \(U\) is open, then so is \(\inv{(g \circ f)}U\).\footnote{It is a
    fact of Set Theory that for \(X\), \(Y\) and \(Z\) sets and
    \(f : X \to Y\) and \(g : Y \to Z\) functions, we have
    \(\inv{(g \circ f)}U = \inv f \left(\inv g U\right)\) for every
    \(U \subseteq Z\).} We can state the following basic properties for the
  composition of continuous functions:
  \begin{tcbenum}
  \item Every topological space \(A\) has associated the continuous
    function
    \[\id_A : A \to A \,,\ \id_A (x) := x\]
    is such that for every topological space \(B\) and continuous
    function \(g : B \to A\) we have
    \[\id_A \circ g = g\]
    and for every topological space \(C\) and continuous function
    \(h : A \to C\) we have
    \[h \circ \id_A = h .\]
  \item \(\circ\) is {\em associative}, that is for \(A\), \(B\), \(C\) and
    \(D\) topological spaces and
    \[A \functo{f} B \functo{g} C \functo{h} D\] continuous functions,
    we have
    \[(h \circ g) \circ f = h \circ (g \circ f) .\]
  \end{tcbenum}
  % Being continuous functions functions, the associativity comes for
  % free; moreover, the identity functions are continuous.
  That is: take the properties listed in the previous example and
  replace \q{set} with \q{topological space} and \q{function} with
  \q{continuous function}.
\end{example}

% \begin{exercise}
%   In Measure Theory, we have \(\sigma\)-{\em algebras}, that is sets where
%   some of its subsets are said to be {\em measurable}. We can define
%   {\em measurable functions} too, that is functions that care about
%   the property od being measurable as continuous functions do of the
%   property of being open.\footnote{Perhaps, you are taught that
%   measurable functions are functions \(f : \Omega \to \reals\) from a
%   measurable space \(\Omega\) such that \(\inv f (-\infty, a]\) is a measurable
%   for every \(a \in \reals\). Anyway, \(\reals\) has the Borel
%   \(\sigma\)-algebra, which is defined as the smallest of the
%   \(\sigma\)-algebras containing the open subsets of \(\reals\) under the
%   Euclidean topology. It can be easily shown that
%   \(f : \Omega \to \reals\) is measurable if and only if \(\inv f B\) is
%   measurable for every Borel subset \(B\) of \(\reals\).} Of course,
%   you can found other example of categories in Algebra, Linear
%   Algebra, Geometry and Analysis. Go and catch as many as you can
%   within you mathematical knowledge. And yes, it may be boring
%   sometimes, and you are right, but as we progress there are
%   remarkable differences from one category to another one.
% \end{exercise}

It should be clear a this point what the pattern is:

% \begin{definition}[Categories]
%   A {\em category} amounts at assigning some things called {\em
%   objects} and, for each couple of objects \(a\) and \(b\), other
%   things named {\em morphisms} from \(a\) to \(b\). We write
%   \(f : a \to b\) to say that \(f\) is a morphism from \(a\) to \(b\),
%   where \(a\) is the {\em domain} of \(f\) and \(b\) the {\em
%   codomain}. Besides, for \(a\), \(b\) and \(c\) objects and
%   \(f : a \to b\) and \(g : b \to c\) morphisms, there is associated the
%   {\em composite morphism}
%   \[gf : a \to c .\] All those things are regulated by the following
%   axioms:
%   \begin{tcbenum}
%   \item for every object \(x\) there is a morphism, \(\id_x\), from
%     \(x\) to \(x\) such that for every object \(y\) and morphism
%     \(g : y \to x\) we have
%     \[\id_x g = g\]
%     and for every object \(z\) and morphism \(h : x \to z\) we have
%     \[h \id_x = h ;\]
%   \item for \(a\), \(b\), \(c\) and \(d\) objects and morphisms
%     \[a \functo{f} b \functo{g} c \functo{h} d\] we have the identity
%     \[(h g) f = h (g f) .\]
%   \end{tcbenum}
% \end{definition}

\begin{definition}[Categories]\label{definition:categories}
  A {\em category} \(\mathcal{C}\) consists of:
  \begin{tcbitem}
  \item A class of {\em objects}, denoted \(\obj{\mathcal{C}}\).
  \item For every \(a, b \in \obj{\mathcal{C}}\), a class of {\em morphisms}
    (or {\em maps} or {\em arrows}) from \(a\) to \(b\), written as
    \(\mathcal{C}(a, b)\). Often, when the category we are working in is
    clear from the context, we write \(f : a \to b\) in place of
    \(f \in \mathcal{C}(a, b)\).
  \item For every \(a, b, c \in \obj{\mathcal{C}}\), an operation of {\em
      composition}
    \[
      \circ_{a, b, c} : \mathcal{C}(b, c) \times \mathcal{C}(a, b) \to \mathcal{C}(a, c) \,, \ (g,
      f) \to g \circ_{a, b, c} f
    \]
    It is customary to write \(g \circ f\) instead of
    \(g \circ_{a, b, c} f\) because the latter can be uncomfortable and
    \(a, b, c\) can be easily inferred from \(f\) and \(g\).
  \end{tcbitem}
  with the following axioms:
  \begin{tcbenum}
  \item For every \(x \in \obj{\mathcal{C}}\) there exists a
    \(\id_x \in \mathcal{C}(x, x)\) such that for every
    \(y \in \obj{\mathcal{C}}\) and \(g \in \mathcal{C}(y, x)\) we have
    \[
      \id_x g = g
    \]
    and for every \(z \in \obj{\mathcal{C}}\) and \(h \in \mathcal{C}(x, z)\) we have
    \[
      h \id_x = h .
    \]
  \item For \(a, b, c, d \in \obj{\mathcal{C}}\) and \(f \in \mathcal{C}(a, b)\),
    \(g \in \mathcal{C}(b, c)\) and \(h \in \mathcal{C}(c, d)\) we have the identity
    \[
      (h g) f = h (g f) .
    \]
  \end{tcbenum}
\end{definition}

% The category of Example~\ref{example:Set} is called \(\Set\) and the
% category of Example~\ref{example:Top} is written \(\Top\).
We list some famous categories here and leave the readers to verify
categorial axioms by themselves.

\begin{center}
  \begin{tabular}{l|ll}
    name & objects & morphisms \\
    \midrule
    \(\Set\) & sets & functions \\
    \(\Top\) & topological spaces & continuous functions \\
    \(\Grp\) & groups & group homomorphisms \\
    \(\Mon\) & monoids & monoid homomorphism \\
    \(\Ring\) & rings & ring homomorphisms \\
    % \(\Rng\) & rings & ring homomorphisms \\
    \(\CRing\) & commutative rings & ring homomorphisms
    \\
    \(\Modu_R\) & modules over a ring & linear maps \\
    \(\Vect_k\) & vector spaces over a field \(k\) & linear maps \\
    \(\FDVect_k\) & finite dimensional vector spaces & linear maps \\
  \end{tabular}
  
\end{center}

\begin{remark}
  We have used the word ``class'' in this definition. The fact that
  there is no ``set of all sets'' is folklore nowadays.\footnote{If we
    want a set \(X\) to be the set of all sets, then it has all its
    subsets as elements, which is an absurd. In fact, Cantor's Theorem
    states that for every set \(X\) there is no surjective function
    \(f : X \to 2^X\).}  So if we want a category like \(\Set\), we need
  classes. All the categories in the table above are too big for
  sets. Since we do not want to exclude such interesting categories, we
  implicitly deploy the Set Theory of von Neumann, Bernays and
  G\"odel. There are classes: they can be either {\em proper classes} or
  {\em sets}.
\end{remark}

We have started with sets and functions, afterwards we have made an
example based on the previous one; if you have accepted the invite of
the exercise above, you have likely found categories where objects are
sets at their core and morphisms are functions with extra property. We
have given an abstract definition of categories because out there are
many other categories that deserve attention.

\begin{example}[Monoids are categories]
  Consider a category \(\mathcal{G}\) with a single object, that we indicate
  with a bare \(\bullet\). All of its morphisms have \(\bullet\) as domain and
  codomain: then any two morphisms are composable and the composite of
  two morphisms \(\bullet \to \bullet\) is a morphism
  \(\bullet \to \bullet\). This motivates us to proceed as follows: let
  \(G\) be the collection of the morphisms of \(\mathcal{G}\) and consider
  the operation of composing morphism
  \[G \times G \to G\,,\ (x, y) \to xy .\] Being \(\mathcal{G}\) a category implies
  this function is associative and \(\mathcal{G}\) has the identity of
  \(\bullet\), that is \(G\) has one element we call \(1\) and such that
  \(f1 = 1f = f\) for every \(f \in G\). In other words, we are saying
  \(G\) is a monoid. We say the single object category \(\mathcal{G}\) {\em
    is} a monoid. \newline Conversely, take a monoid \(G\) and a symbol
  \(\bullet\): make such thing acquire the status of object and the elements
  of \(G\) that of morphisms; in that case, the operation of \(G\) has
  the right to be called composition because the axioms of monoid say
  so. Here, \(\bullet\) is something we care of just because by definition
  morphisms require objects and it has no role other than this.
\end{example}

In Mathematics, a lot of things are monoids, so this is nice. In
particular, a {\em group} is a single object category where for every
morphism \(f\) there is a morphism \(g\) such that \(gf\) and \(fg\) are
the identity of the the unique object. We will deal with isomorphism
later in this chapter.

\begin{example}[Preordered sets are categories]
  A {\em preordered set} (sometimes contracted as {\em proset}) consists
  of a set \(A\) and a relation \(\le\) on \(A\) such that:
  \begin{tcbenum}
  \item \(x \le x\) for every \(x \in A\);
  \item for every \(x, y, z \in A\) we have that if \(x \le y\) and
    \(y \le z\) then \(x \le z\).
  \end{tcbenum}
  Now we do this: for \(x, y \in A\), whenever \(x \le y\) take
  \((a, b) \in A \times A\). We operate with these couples as follows:
  \begin{equation}
    (y, z) (x, y) := (x, z), \label{eqn:ProsetCirc}
  \end{equation}
  where \(x, y, z \in A\). This definition is perfectly motivated by (2):
  in fact, if \(x \le y\) and \(y \le z\) then \(x \le z\), and so there is
  \((x, z)\). By (1), for every \(x \in A\) we have the couple
  \((x, x)\), which has the following property: for every \(y \in A\)
  \begin{equation}
    \begin{aligned}
      (x, y) (x, x) &= (x, y) & \text{for every } y \in A \\
      (x, x) (z, x) &= (z, x) & \text{for every } z \in A .
    \end{aligned}\label{eqn:ProsetId}
  \end{equation}
  Another remarkable feature is that for every
  \(x_1, x_2, x_3, x_4 \in A\)
  \begin{equation}
    \big((x_3, x_4)(x_2, x_3)\big)(x_1, x_2) = (x_3, x_4)\big((x_2,
    x_3)(x_1, x_2)\big)
    \label{eqn:ProsetAssoc}
  \end{equation}
  We have a category indeed: its objects are the elements of \(A\), the
  morphisms are the couples \((x, y)\) such that \(x \le y\)
  and~\eqref{eqn:ProsetCirc} gives the notion of composition;
  \eqref{eqn:ProsetId} says what are identities
  while~\eqref{eqn:ProsetAssoc} tells the compositions are associative.
\end{example}

Several things are prosets, so this is nice. Namely, {\em partially
  ordered sets}, or {\em posets}, are prosets where every time there are
morphisms going opposite directions
\[\begin{tikzcd}
    a \ar[r, shift left] & b \ar[l, shift left]
  \end{tikzcd}\] then \(a = b\). Later in this chapter, we will meet
{\em skeletal} categories.

\begin{example}[Matrices]\label{example:Mat}
  Fixed some field \(k\), for \(m\) and \(n\) positive integers, a {\em
    matrix} of type \(m \times n\) is just an arrangement of elements of
  \(k\) in \(m\) rows and \(n\) columns as follows:
  \[
    \begin{pmatrix}
      x_{1,1} & x_{1,2} & \cdots{} & x_{1,n} \\
      x_{2,1} & x_{2,2} & \cdots{} & x_{2,n} \\
      \vdots  & \vdots  &          & \vdots  \\
      x_{m,1} & x_{m,2} & \cdots{} & x_{m,n}
    \end{pmatrix}
  \]
  If \(A\) is the name of a matrix, then \(A_{i, j}\) is the element on
  the intersection of the \(i\)th row and the \(j\)th column. Matrices
  can be multiplied: if \(A\) and \(B\) are matrices of type
  \(m \times n\) and \(n \times r\) respectively, then \(AB\) is the matrix of
  type \(m \times r\) where
  \[
    (AB)_{i, j} := \sum_{p = 1}^n A_{i, p} B_{p ,j} .
  \]
  Our experiment is the following. Consider the positive integers in the
  role of objects and, for \(m\) and \(n\) integers, the matrices of
  type \(m \times n\) as morphisms from \(n\) to \(m\). Further, if
  \(A\) ha type \(m \times n\) and \(B\) has type \(n \times r\), define the
  composition of
  \[
    \begin{tikzcd}[column sep=small]
      r \ar["{B}", r] & n \ar["{A}", r] & m
    \end{tikzcd}
  \]
  as the matrix product \(AB\). Let us investigate whether categorial
  axioms hold.
  \begin{tcbitem}
  \item For \(n\) positive integer, we have the {\em identity matrix}
    \(I_n\), the one of type \(n \times n\) defined by
    \[
      (I_n)_{i, j} =
      \begin{cases}
        1 & \text{if } i = j \\
        0 & \text{otherwise}
      \end{cases}
    \]
    One, in fact, can verify that such matrix is an identity in
    categorial sense: for every positive integer \(m\), an object, and
    for every matrix \(A\) of type \(m \times n\), a morphism from \(n\) to
    \(m\), we have
    \[
      A I_n = A
    \]
    and for every positive integer \(r\) and for every matrix \(B\) of
    type \(r \times n\) we have
    \[
      I_n B = B .
    \]
  \item For \(A\), \(B\) and \(C\) matrices of type \(m \times n\),
    \(n \times r\) and \(r \times s\) respectively, we have
    \[
      (AB)C = A(BC) .
    \]
    Again, this identity can be regarded under a categorial light.
  \end{tcbitem}
  The category of matrices over a field \(k\) just depicted is written
  \(\Mat_k\). This example may seem quite useless, but it really does
  matter when you know there is the category of finite vector spaces
  \(\FDVect_k\): just wait until we talk about equivalence of
  categories.
\end{example}



\section{Functors}

\begin{definition}[Functors]\label{definition:functors}
  A {\em functor} \(F\) from a category \(\mathcal{C}\) to a category
  \(\mathcal{D}\) amounts to:
  \begin{tcbitem}
  \item One \q{function on objects}
    \[
      F_{\text{obj}} : \obj{\mathcal{C}} \to \obj{\mathcal{D}}
    \]
  \item For every \(a, b \in \obj{\mathcal{C}}\), one \q{function on morphisms}
    \[
      F_{a, b} : \mathcal{C}(a, b) \to \mathcal{D}(F_{\text{obj}}(a),
      F_{\text{obj}}(b))
    \]
  \end{tcbitem}
  such that:
  \begin{tcbenum}
  \item For every \(x \in \obj{\mathcal{C}}\), we have
    \(F_{x, x}(\id_x) = \id_{F_{\text{obj}}(x)}\).
  \item For every \(x, y, z \in \obj{\mathcal{C}}\) and \(f : x \to y\) and
    \(g : y \to z\) in \(\mathcal{C}\), we have
    \(F_{b, c}(g) F_{a, b}(f) = F_{a, c}(gf)\).
  \end{tcbenum}
  To say that \(F\) is a functor from \(\mathcal{C}\) to \(\mathcal{D}\) we use
  \(F : \mathcal{C} \to \mathcal{D}\), a symbolism that recalls that one of morphism
  in categories. We write \(F(a)\) instead of \(F_{\text{obj}}(a)\) if
  \(a\) is an object and \(F(f)\) instead of \(F_{a, b}(f)\) if
  \(f : a \to b\) because there cannot be any confusion if the context is
  clear.
\end{definition}

A first straightforward consequence of functoriality is contained in the
following proposition.

\begin{proposition}
  Let \(F : \mathcal{C} \to \mathcal{D}\) be a functor. If \(f\) is an isomorphism
  of \(\mathcal{C}\), then so is \(F(f)\).
\end{proposition}

\begin{proof}
  \inlinethm{Exercise}.
\end{proof}

As often happens, let us start with simple exmaples: in this context,
the simplest ones can be obtained by choosing very simple categories.

\begin{example}[Functors from sets]\label{example:CollectionsAreCats}
  You can identify a set to the category where its objects are the
  elements of the class and the unique morphisms present are the
  identities. Consider a functor \(F : \mathcal{S} \to \mathcal{C}\) where \(S\) is
  a set. \(F\) maps objects to objects and morphisms to morphisms, but
  the only morphisms of \(\mathcal{S}\) are identities, which are taken to
  identities of \(\mathcal{C}\). Since \(F\) involves only objects and
  identities, \(F\) is just a function \(\obj{\mathcal{S}} \to \obj{\mathcal{C}}\).
\end{example}

\begin{example}[Functors from prosets]
  Consider a functor \(F : (A, \le) \to \mathcal{C}\) from a proset. We know
  \((A, \le)\) is a category that has at most one morphism for each
  ordered couple in \(A \times A\). For that reason, let us adopt this
  notation: for every \(i, j \in A\) such that \(i \le j\) indicate by
  \(F_{i,j}\) the image of the unique morphism \(i \to j\) of
  \((A, \le)\) via \(F\). That being said, our functor \(F\) is just a
  collection \(\set{F_i \mid i \in A}\) with morphisms \(F_{i,j}\) for
  \(i, j \in A\) if \(i \le j\).
  % As a particular instance of this, let us
  % examine functors
  % \[H : (\naturals, \le) \to \mathcal{C}\] with \(\le\) being the usual ordering of
  % \(\naturals\). For \(i, j \in \naturals\) with \(i \le j\), the morphism
  % \(i \to j\) can be factored into consecutive morphisms
  % \[i \to j = (j-1 \to j) \cdots{} (i+1 \to i+2) (i \to i+1) .\] For that reason, our
  % \(H\) \q{is} just a sequence
  % \[H_0 \functo{\partial_0} H_1 \functo{\partial_1} \cdots{} \functo{\partial_{n+1}} H_n
  %   \functo{\partial_n} \cdots{}\] of objects and morphisms in
  % \(\mathcal{C}\), where we have written \(\partial_j\) for \(H_{j, j+1}\).
\end{example}

\begin{example}[Monotonic functions]
  Consider two prosets \((A, \le_A)\) and \((B, \le_B)\). A functor
  \(F : (A, \le_A) \to (B, \le_B)\) maps objects to objects and morphisms to
  morphisms in the following way: for every \(a, b \in A\) if
  \(a \le_A b\) then \(F(a) \le_b F(b)\). In other words, \(F\) is a {\em
    monotonic function}.
\end{example}

\begin{example}[Monoid homomorphsisms]
  We have previously seen that monoids are single object
  categories. Consider now two such categories, say \(\mathcal{G}\) and
  \(\mathcal{H}\), and let us see what a functor \(f : \mathcal{G} \to \mathcal{H}\)
  is. Denoting by \(\bullet_{\mathcal{G}}\) and \(\bullet_{\mathcal{H}}\) the object of
  \(\mathcal{G}\) and \(\mathcal{H}\) respectively, there is a unique possibility:
  mapping \(\bullet_{\mathcal{G}}\) to \(\bullet_{\mathcal{H}}\). The functorial axioms in
  that case are:
  \[
    f(xy) = f(x)f(y)
  \]
  for every morphisms \(x\) and \(y\) of \(\mathcal{G}\) and
  \[
    f(1_{\mathcal{G}}) = 1_{\mathcal{H}} ,
  \]
  with \(1_{\mathcal{G}}\) and \(1_{\mathcal{H}}\) being the identities of
  \(\mathcal{G}\) and \(\mathcal{H}\) respectively. These two properties say that
  \(f\) is a monoid homomorphism. In this case there is also an equation
  that about objects, but these two are a mere subtlety that adds
  nothing.
\end{example}

Let us remain to Algebra. What happens if the domain of a functor is a
group and the codomain is \(\Set\)? (Recall that \q{group} is just a
fancier abbreviation of \q{single object groupoid}.) The answer is
contained in the following example.

\begin{example}[Group actions]
  If \(\mathcal{G}\) is a group, what is a functor
  \(\delta : \mathcal{G} \to \Set\)? Let us write \(G\) the set of its
  morphisms. The unique object of \(\mathcal{G}\) is mapped to one set
  \(X\). Since all the morphisms of \(\mathcal{G}\) are isomorphisms, then
  \(\delta\) takes each of them to bijections from \(X\) to \(X\). Thus we
  can say \(\delta\) is the assignment of a certain set \(X\) and, for every
  \(g\) element of the group, of one isomorphism \(\delta_g : X \to X\), done
  in a very specific manner:
  \begin{itemize}
  \item \(\delta_{\id} = \id_X\), where \(\id\) is the identity on the unique
    object of \(\mathcal{G}\).
  \item \(\delta_{gf} = \delta_g \delta_f\) for every \(f, g \in G\).
  \end{itemize}
  Does this sound familiar? What we have described is a group {\em
    action} over the set \(X\), that is a group isomorphism from \(G\)
  to the symmetric group of \(X\).
\end{example}

\begin{figure}
  \centering
  \input{./tikz/group-actions.qtikz}
  \caption{A group action as functor}
\end{figure}

\begin{example}[The category \(\Eqv\)]\label{example:CategoryEqv}
  A {\em setoid} \NotaInterna{nlab uses this term\dots{}} is having a
  set and an equivalence relation defined on it. If \(X\) is a set and
  \(\sim\) an equivalence relation over \(X\), the setoid amounting of
  these data is written as \((X, \sim)\). Any set \(X\) has of course its
  own equality, that we denote by \(=_X\).\footnote{In Set Theory,
    \(=_X\) is the set \(\set{(a, a) \mid a \in X}\).}\newline For if
  \(X\) and \(Y\) sets, a function \(f : X \to Y\) respects this rule by
  definition:
  \begin{quotation}
    for every \(a, b \in X\), if \(a =_X b\) then \(f(a) =_Y b\).
  \end{quotation}
  We would like to replace the equalities above with equivalence
  relations: for if \((X, \sim_X)\) and \((Y, \sim_Y)\) are setoids, a {\em
    functoid} \NotaInterna{ok, let me find/craft a nicer name\dots{}}
  from \((X, \sim_X)\) to \((Y, \sim_Y)\) is any function
  \(f : X \to Y\) such that
  \begin{quotation}
    for every \(a, b \in X\), if \(a \sim_X b\) then \(f(a) \sim_Y f(b)\).
  \end{quotation}
  Functoids are certain type of functions, and composing two of them as
  such returns a funtoid. Categorial axioms hold almost for free, so we
  really have a {\em category of setoids and functoids},
  \(\Eqv\).\newline Let us now involve functoriality. There is a nice theorem:
  \begin{quotation}
    Let \(X\) and \(Y\) be two sets with \(\sim_X\) and \(\sim_Y\) equivalence
    relations on \(X\) and \(Y\) respectively. Then for every
    \(f : X \to Y\) such that \(f(a) \sim_Y f(b)\) for every
    \(a, b \in X\) such that \(a \sim_X b\), there exists one and only one
    \(\phi : X{/}{\sim_X} \to Y{/}{\sim_Y}\) that makes
    \[
      \begin{tikzcd}[column sep=small]
        X \ar["f", r] \ar["{\lambda a.[a]_X}", d, swap] & Y \ar["{\lambda b.[b]_Y}", d] \\
        X{/}{\sim_X} \ar["\phi", r, swap] & Y{/}{\sim_Y}
      \end{tikzcd}
    \]
    commute. (The vertical functions are the canonical projections.)
  \end{quotation}
  This underpins the functor
  \[
    \pi : \Eqv \to \Set
  \]
  that maps setoids \((X, \sim)\) to sets \(X{/}{\sim}\) and functoids
  \(f : (X, \sim_X) \to (Y, \sim_Y)\) to functions
  \[
    \begin{aligned}
      & \pi_f : X{/}{\sim_X} \to Y{/}{\sim_Y} \\
      & \pi_f \left([a]_X\right) := \left[f(a)\right]_Y ,
    \end{aligned}
  \]
  whose existence and uniqueness is claimed by the just mentioned
  Proposition.
\end{example}

\begin{example}[Free groups]\label{example:FreeGroup}
  Suppose given a {\em group alphabet} \(S\), which is just a set of
  things we decide to name \q{letters}. Then a {\em group word} on \(S\)
  is a string obtained by juxtaposition of a finite amount of
  \q{\(x^1\)} and \q{\(x^{-1}\)}, where \(x \in S\). The {\em empty word}
  is obtained by writing no letter at all, and we shall denote it by
  something, say \(e\); instead, the other words appear as
  \[
    x_1^{\phi_1} \cdots{} x_n^{\phi_n}
  \]
  with \(x_1, \dots{}, x_n \in S\) and
  \(\phi_1, \dots{}, \phi_n \in \set{-1, 1}\). We shall add another constraint
  on what is a word: for \(i \in \set{1, \dots{}, n-1}\) if
  \(x_i = x_{i+1}\), then \(\phi_i \ne - \phi_{i+1}\). In other words,
  substrings of the form \(xx^{-1}\) and \(x^{-1}x\) are not
  allowed. Let us write \(\angled S\) the set of all words written using
  the alphabet \(S\).\newline The length of a word is the number of letters it
  is made of. We define equality only on words having the same length:
  we say \(x_1^{\alpha_1} \cdots{} x_n^{\alpha_n}\) is equal to
  \(y_1^{\beta_1} \cdots{} y_n^{\beta_n}\) whenever \(x_i = y_i\) and
  \(\alpha_i = \beta_i\) for every
  \(i \in \set{1, \dots{}, n}\).\newline It is natural to join two words by bare
  juxtaposition, but the resulting word may not be irreducible; this
  issue has to be fixed:
  \begin{align*}
    & \cdot : \angled S \times \angled S \to \angled S \\
    & e \cdot w := w \,, \ w \cdot e := w \\
    & (x_1^{\lambda_1} \cdots x_m^{\lambda_m}) \cdot (y_1^{\mu_1} \cdots y_n^{\mu_n}) :=
      \begin{cases}
        (x_1^{\lambda_1} \cdots x_{m-1}^{\lambda_{m-1}}) \cdot (y_2^{\mu_2} \cdots y_n^{\mu_n}) & \text{if } x_{m}^{\lambda_m} = y_1^{-\mu_1} \\
        x_1^{\lambda_1} \cdots x_m^{\lambda_m} y_1^{\mu_1} \cdots y_n^{\mu_n} & \text{otherwise.}
      \end{cases}
  \end{align*}
  Let us define a function that either reverses the order of the letters
  and changes each exponent to the other one:
  \[
    i : \angled S \to \angled S \,, \ i\left( x_1^{\xi_1} \cdots{} x_i^{\xi_i}
      x_{i+1}^{\xi_{i+1}} \cdots{} x_n^{\xi_n} \right) := x_n^{-\xi_n} \cdots{}
    x_{i+1}^{-\xi_{i+1}} x_i^{-\xi_i} \cdots{} x_1^{-\xi_1} .
  \]
  It is immediate to show that \(w \cdot i(w) = i(w) \cdot x = e\) for every
  \(w \in \angled S\). Only the associativity of \(\cdot\) is a a bit tricky
  to prove. At this point we have endowed \(\angled S\) with a group
  structure.\newline Now, if take two sets \(S\) and \(T\) and a function
  \(f : S \to T\), we have the group homomorphism
  \begin{align*}
    & \angled f : \angled S \to \angled T \\
    & \angled f (x_1^{\delta_1} \cdots x_n^{\delta_n}) := f(x_1)^{\delta_i} \cdots f(x_n)^{\delta_n} .
  \end{align*}
  It is immediate to demonstrate that we ended up with having a functor
  \[
    \angled{\phantom\square} : \Set \to \Grp .
  \]
\end{example}

% In future we will provide other examples involving free modules.

% \begin{example}[Free modules]
%   The explicit construction of the free {\em abelian} group given a set
%   is simpler than that of free group in general. Since an abelian group
%   is a \(\integers\)-module, let us show how to build a free module.\newline
%   Let \(R\) be a ring and \(S\) be a set, as in the previous
%   example. Intuitively, the module generated by \(S\) are linear
%   combination of a finite amount of elements of \(S\), that is
%   expressions of the form
%   \[\sum_{i = 1}^n \lambda_i x_i\]
%   for \(n \in \naturals\), \(\lambda_1, \dots{}, \lambda_n \in R\) and
%   \(x_1, \dots{}, x_n \in S\). Observe, however that the \q{sum} here is
%   just a formal expression: there is no link to the an operation of sum
%   yet. We can rethink this linear combination as something more
%   manageable during computations:
%   \[\sum_{x \in S} \lambda_x x\]
%   where \(\lambda : S \to R\) is non zero for a finite amount of elements of
%   \(S\). Observe that it is just formalism: \(S\) may be an infinite
%   set, but the sum \(\sum_{x \in S} \lambda_x x\) is not to be understood as a
%   series in Analysis; consider also \(\lambda_x \ne 0\) for finitely many
%   \(x\), so if \(S\) is infinite, the most of the terms are
%   useless. %\newline
%   \NotaInterna{Instead of using the device of \q{formal expressions}, we
%     can define the module words as functions \(\lambda : S \to R\) that assume
%     non-zero values for a finite amount of elements. Isn't that the same
%     stuff of a formal sum?}\newline Thus, let us write the explicit definition:
%   \[
%     \angled S := \set{\left.\sum_{x \in S} \lambda_x x \right\mid \lambda : S \to R\,,\, \lambda_x \ne
%       0 \text{ for finitely many } x \in S} .
%   \]
%   This is only the first step to make a module with such set. We give a
%   sum
%   \begin{align*}
%     & + : \angled S \times \angled S \to \angled S \\
%     & \left( \sum_{x \in S} \alpha_x x \right) + \left( \sum_{x \in S} \beta_x x \right) : = \sum_{x \in S} (\alpha_x + \beta_x) x
%   \end{align*}
%   and an external product
%   \begin{align*}
%     & \cdot : R \times \angled S \to \angled S \\
%     & \eta \cdot \left( \sum_{x \in S} \alpha_x x \right) : = \sum_{x \in S} (\eta \alpha_x) x
%   \end{align*}
%   It is simple to verify that \(\angled S\) is a \(R\)-module now.\newline So
%   far, we only have an process that takes sets and emits \(R\)-modules:
%   to make a functor, we also need to instruct how to construct a linear
%   function from a simple function of sets. For \(f : S \to T\), we give
%   \begin{align*}
%     & \angled f : \angled S \to \angled T \\
%     & \angled f \left(\sum_{x \in S} \lambda_x x \right) := \sum_{x \in S} \lambda_x f(x) .
%   \end{align*}
%   It is simple to verify we have a functor
%   \[\angled\hole : \Set \to \Modu_R .\]
% \end{example}

\begin{exercise}
  There is a plenty of \q{free stuff} around that can give arise to
  functors like the one above. Find and illustrate some of them.
\end{exercise}

\begin{example}[The First Homotopy Group]
  A {\em pointed topological space} is a topological space \(X\) with
  one point \(x_0 \in X\); we write it as \((X, x_0)\). We define a {\em
    pointed continuous function} \((X, x_0) \to (Y, y_0)\) to be a
  continuous functions \(X \to Y\) taking \(x_0\) to \(y_0\). Furthermore,
  composing such functions yields a pointed continuous function. So,
  really we have the category of pointed topological spaces, we denote
  by \(\Top_\ast\).
  \[
    \Omega(X, x_0) := \set{\text{continuous } \phi : [0,1] \to X \mid \phi(0) = \phi(1) =
      x_0}
  \]
  and call its elements {\em loops} of \(X\) based at \(x_0\).  Two
  loops can be joined, that is traversing one loop after another one:
  for if \(\phi, \psi \in \Omega(X, x_0)\) we introduce the loop
  \(\phi \ast \psi : [0,1] \to X\) with
  \[
    (\phi \ast \psi) (t) :=
    \begin{cases}
      \phi(2t) & \text{if } t \le \frac12 \\
      \psi(2t-1) & \text{otherwise.}
    \end{cases}
  \]
  This gives us the operation of {\em junction} of loops
  \[
    \ast : \Omega(X, x_0) \times \Omega(X, x_0) \to \Omega(X, x_0) .
  \]
  Now it's time to find a suitable equivalence relation that is
  compatible with this operation. For if
  \(\phi, \psi \in \Omega(X, x_0)\), we say \(\phi\) is {\em homotopic} to
  \(\psi\) whenever there exists a {\em homotopy} from \(\phi\) to
  \(\psi\), viz a continuous function
  \[H : [0,1] \times [0,1] \to X\] such that
  \(H(\phantom{s}, 0) = \phi\), \(H(\phantom{s}, 1) = \psi\),
  \(H(s, 0) = H(s, 1) = x_0\) for every \(s \in [0, 1]\). This relation is
  an equivalence one and it is compatible with \(\ast\). It remains to
  verify some properties to define a group structure:
  \begin{tcbitem}
  \item \((\alpha \ast \beta) \ast \gamma\) is homotopic to
    \(\alpha \ast (\beta \ast \gamma)\) for every \(\alpha, \beta, \gamma \in \Omega(X, x_0)\).
  \item the paths \(\alpha \ast c_{x_0}\), \(c_{x_0} \ast \alpha\) and
    \(\alpha\) are homotpic for every \(\alpha \in \Omega(X, x_0)\); here,
    \(c_{x_0}\) is the loop defined by \(c_{x_0}(t) = x_0\).
  \item \(\alpha \ast \inv\alpha\), \(\inv \alpha \ast \alpha\) and
    \(c_{x_0}\) are homotopic for every \(\alpha \in \Omega(X, x_0)\); here,
    \(\inv \alpha\) is the loop with \(\inv \alpha (t) := \alpha (1-t)\).
  \end{tcbitem}
  We have now all the ingredients to introduce a group: define
  \(\pi_1(X, x_0)\) to be the set obtained identifying homotopic elements
  of \(\Omega(X, x_0)\); this set is a group once you consider the operation
  \begin{align*}
    \pi_1(X, x_0) \times \pi_1(X, x_0) &\to \pi_1(X, x_0) \\
    ([\alpha], [\beta]) &\to [\alpha][\beta] := [\alpha \ast \beta] .
  \end{align*}
  Here, we have written \([\phi]\) for the set of loops homotopic to
  \(\phi\). Sometimes --- especially if we are considering more topological
  spaces ---, we need to specify the topological space we are taking
  loops, for example writing \([\phi]_X\). Now, it is the turn to define
  induced homomorphisms: for a pointed continuous function
  \(f : (X, x_0) \to (Y, y_0)\) we have the group homomorphism
  \begin{align*}
    &\pi_1 (f) : \pi_1(X, x_0) \to \pi_1(Y, y_0) \\
    & \pi_1(f)[\phi]_X := [f \phi]_Y .
  \end{align*}
  (You may have a look at Example~\ref{example:CategoryEqv}.)  Instead
  of \(\pi_1(f)\), you may have been get used to \(f_\ast\). In conclusion,
  we have just defined one functor
  \[
    \pi_1 : \Top_\ast \to \Grp .
  \]
  The {\em first fundamental group} is not just a group, and that the
  actual group is just a piece of larger picture.
\end{example}

\begin{example}[Probability distributions]
  In Probability and Statistics we have a construction that occurs
  almost everywhere:
  \begin{quotation}
    Given a probability space \((\Omega, \mathcal{A},\mathbb{P})\) a measurable space
    \((\Gamma, \mathcal{B})\) and a measurable function \(X : \Omega \to \Gamma\), then
    \[
      \mathbb{P}_X : \mathcal{B} \to [0,1] \,,\ \mathbb{P}_X(E) := \mathbb{P}\left(X^{-1}E\right) = \mathbb{P}[X \in E]
    \]
    is a probability measure on \((\Gamma, \mathcal{B})\).
  \end{quotation}
  How do we put categories here? Well, let us rephrase the theorem
  above: if we drop the probability from the probability space we have
  the measurable space \((\Omega, \mathcal{A})\); there is also the other measurable
  space and the measurable function; if you think closely, under this
  context the theorem allows to construct from a generic probability
  measure on \((\Omega, \mathcal{A})\) a probability measure on
  \((\Gamma, \mathcal{B})\). So here we go.
  \begin{quotation}
    Given a measurable function
    \(X : (\Omega, \mathcal{A}) \to (\Gamma, \mathcal{B})\), there is the function
    \begin{align*}
      \set{\text{measures on } (\Omega, \mathcal{A})} & \to \set{\text{measures on }
                                         (\Gamma, \mathcal{B})} \\
      \mathbb{P} & \mapsto \mathbb{P}_X
    \end{align*}
  \end{quotation}
  The measurable function
  \(X : (\Omega, \mathcal{A}) \to (\Gamma, \mathcal{B})\) lives in the category of measurable spaces
  \(\Meas\) whereas the other function lives in \(\Set\).

  Well, seems a situation where functors are involved. Define
  \(\operatorname{Prob}(\Omega, \mathcal{A})\) to be the set of probability measures on
  \((\Omega, \mathcal{A})\) and
  \[
    \operatorname{Prob} \left(
      \begin{tikzcd}
        (\Omega, \mathcal{A}) \ar["X", d] \\ (\Gamma, \mathcal{B})
      \end{tikzcd}
    \right)
  \]
  as the function
  \(\operatorname{Prob}(\Omega, \mathcal{A}) \to \operatorname{Prob}(\Gamma, \mathcal{B})\) that takes
  \(\mathbb{P}\) to \(\mathbb{P}_X\), as explained above. You might have seen written
  \(X_\ast\) instead of \(\operatorname{Prob}(X)\).

  We check now that we have a functor
  \[
    \operatorname{Prob} : \Meas \to \Set
  \]
\end{example}

Functors can be composed --- and I think at this point it is not a
secret. Take \(\mathcal{C}\), \(\mathcal{D}\) and \(\mathcal{E}\) categories and
functors
\[\mathcal{C} \functo F \mathcal{D} \functo G \mathcal{E} .\]
The sensible way to define the composite functor
\(GF : \mathcal{C} \to \mathcal{E}\) is mapping the objects \(x\) of \(\mathcal{C}\) to
the objects \(GF(x)\) of \(\mathcal{E}\), and the morphisms \(f : x \to y\) of
\(\mathcal{C}\) to the morphisms \(GF(f) : GF(x) \to GF(y)\) of \(\mathcal{E}\). That being set, the composition is associative and there is an
identity functor too.

There are all the conditions, so what prevents us to consider a category
--- we can call \(\bf Cat\) --- that has categories as objects and functors
as morphisms?

If we work upon NBG, we can think of any proper class as a category, for
this statement have a closer look at
Example~\ref{example:CollectionsAreCats}. What happens now is that the
class of objects of \(\bf Cat\) has an element that is a proper class,
which isn't legal in NBG.

Is a category of {\em locally small} categories and functors
problematic? Take \(\mathcal{C}\) such that \(\mathcal{C}(a, b)\) is a proper
class for some \(a\) and \(b\) objects: consider \(\mathcal{C}{/}{b}\). In
this case \(\obj{\mathcal{C}/b}\) is a proper class too, and here we go
again.

Now what? If we stick to NBG, this is a limit we have to take into
account. From now on, \(\bf Cat\) is the category of {\em small}
categories and functors between small categories.



\section{The language of diagrams}

A diagram is a drawing made of \q{nodes}, that is empty slots, and
\q{arrows}, that part from some nodes and head to other ones. Here is
an example:
\begin{equation}\begin{tikzcd}[row sep=tiny]
    & \phantom\square \ar[dr, bend left=10] \\
    \phantom\square \ar[ur, bend left=20] \ar[ur, bend right=20, swap] & & \phantom\square  \\
    & \phantom\square \ar[uu, bend right=10, swap]
  \end{tikzcd}\label{diagram:MuteDiag}\end{equation}
%
Nodes are the places where to put objects' names and arrows are to be
labelled with morphisms' names. The next step is putting labels
indeed, something like this:
\begin{equation}
  \begin{tikzcd}[row sep=tiny]
    & b \ar["q", dr, bend left=10] \\
    a \arrow["f", ur, bend left=20] \arrow["g", ur, bend right=20, swap] & & d  \\
    & c \arrow["h", uu, bend right=10, swap]
  \end{tikzcd}
  \label{diagram:LabelDiag}
\end{equation}
%
The idea we want to capture is: having a scheme of nodes and arrows,
as in~\eqref{diagram:MuteDiag}, and then assigning labels, as
in~\eqref{diagram:LabelDiag}. Since diagrams serve to graphically show
some categorial structure, there should exist the possibility to
\q{compose} arrows: two consecutive arrows
\begin{equation}
  \begin{tikzcd}
    \phantom\square \ar[r, bend left=20] & \phantom\square \ar[r, swap, bend
    right=20] & \phantom\square
  \end{tikzcd}
  \label{diagram:ConseqArrs}
\end{equation}
naturally yields that one that goes from the first node and heads to
the last one; if in~\eqref{diagram:ConseqArrs} we label the arrows
with \(f\) and \(g\), respectively, then the composite arrow is to be
labelled with the composite morphism \(gf\). That operation shall be
associative and there should exist identity arrows too, that is arrows
that represent and behave exactly as identity morphisms. In other
words, our drawings shall care of the categorial structure. If we want
to formalise the idea just outlined, functoriality is exactly what we
need:

% Yes, there is a formal definition of diagram, but we'd better defer
% this sophistication a bit later.

\begin{definition}[Diagrams]
  A {\em diagram} in a category \(\mathcal{C}\) consists of a \q{index
    category} \(\mathcal{I}\) and a functor \(X : I \to \mathcal{C}\).
  % A {\em diagram} in a category \(\mathcal{C}\) is having:
  % \begin{tcbitem}
  % \item a scheme of nodes and arrows, that is a category \(\mathcal{I}\);
  % \item labels for nodes, that is for every object \(i\) of
  %   \(\mathcal{I}\) one object \(x_i\) of \(\mathcal{C}\);
  % \item labels for arrows, that is for every pair of objects \(i\)
  %   and \(j\) of \(\mathcal{I}\) and morphism \(\alpha : i \to j\) of
  %   \(\mathcal{I}\), one morphism \(f_\alpha : x_i \to x_j\) of \(\mathcal{C}\)
  % \end{tcbitem}
  % with all this complying the following rules:
  % \begin{tcbenum}
  % \item \(f_{\id_i} = \id_{x_i}\) for every \(i\) object of
  %   \(\mathcal{I}\);
  % \item \(f_\beta f_\alpha = f_{\beta\alpha}\), for \(\alpha\) and
  %   \(\beta\) two consecutive morphisms of \(\mathcal{I}\).
  % \end{tcbenum}
\end{definition}

% Rather than thinking diagrams abstractly --- like in the form stated
% in the definition ---, one usually draws them. In general, it is not a
% good idea to draw all the compositions. For example, consider four
% nodes and three arcs displayed as
% \[
%   \begin{tikzcd}[row sep=tiny]
%     a \ar["f"{description}, dr] & & c \ar["h"{description}, dr] \\
%     & b \ar["g"{description}, ur, swap] & & d
%   \end{tikzcd}
% \]
% and draw all the compositions: you will convince yourself it may be
% a huge mess even for small diagrams. In fact, why waste an arrow to
% represent the composite \(gf\) in
% \[
%   \begin{tikzcd}[row sep=tiny]
%     a \ar["f", dr] & & c \\
%     & b \ar["g", ur, swap]
%   \end{tikzcd}
% \]
% when \(gf\) is walking along \(f\) before and \(g\) then? Neither
% identities need to be drawn: we know every object has one and only
% one identity and thus the presence of an object automatically
% carries the presence of its identity.

Consecutive arrows in a diagram form a \q{path}; in that case, we
refer to the domain of its first arrow as the domain of the path and
to the codomain of the last one as the codomain of the path. Two paths
are said {\em parallel} when they share both domain and codomain. A
diagram is said to be {\em commutative} whenever any pair of parallel
paths yields the same composite morphism.

Well, but do we express this idea in with the definition of diagram we
have crafted? Here are some examples.

\begin{example}[Commutative triangles]
  Consider the category \(\mathbb{3}\) that we draw as follows:
  \[
    \begin{tikzcd}[column sep=tiny]
      & \bullet \ar["j", dr]  \\
      \bullet \ar["i", ur] \ar["k", rr, swap] & & \bullet
    \end{tikzcd}
  \]
  We haven't drawn identities because we assume the are always present
  together the corresponding objects. Here the unique composition is
  defined as follows:
  \[
    ji := k
  \]
  A {\em commutative triangle} in a category \(\mathcal{C}\) is a diagram
  with scheme \(\mathbb{3}\).
\end{example}

\begin{example}[Commutative squares]
  Let \(\mathcal{I}\) be the category
  \[
    \begin{tikzcd}
      \bullet \ar["{i_1}", r] \ar["{j_1}", d, swap] \ar["k", dr] & \bullet \ar["i_2", d] \\
      \bullet \ar["{j_2}", r, swap] & \bullet
    \end{tikzcd}
  \]
  where compositions are defined as follows:
  \begin{align*}
    & i_2i_1 := k \\
    & j_2j_i := k
  \end{align*}
  A {\em commutative square} in a category \(\mathcal{C}\) is a diagram
  whose scheme is \(\mathcal{I}\).
\end{example}

Since diagrams are meant to be drawn and the definition we have given
is just the formalisation of these devices, one says a commutative
triangle is three objects and three morphisms displayed as
\[
  \begin{tikzcd}[column sep=tiny]
    & b \ar["g", dr]  \\
    a \ar["f", ur] \ar["h", rr, swap] & & c
  \end{tikzcd}
\]
and where \(gf = h\). Similarly, one says a commutative square is four
objects and four morphisms displayed as
\[
  \begin{tikzcd}
    a \ar["f", r] \ar["h", d, swap] & c \ar["g", d] \\
    b \ar["k", r, swap] & d
  \end{tikzcd}
\]
where \(gf = kh\). Sometimes, commutative squares are not really
squares. For example
\[
  \begin{tikzcd}
    a \ar["f", r, shift left] \ar["g", r, shift right, swap] & b
    \ar["h", r] & c
  \end{tikzcd}
\]
where \(hf = hg\) is a commutative square as well: the index category is that of the example but
with \(i_1\) and \(j_1\) mapped to \(f\) and \(g\) respectively and
with \(i_2\) and \(j_2\) mapped to \(h\). If you want to a square you
can be repetitive
\[
  \begin{tikzcd}
    a \ar["f", r] \ar["g", d, swap] & b \ar["h", d] \\
    b \ar["h", r, swap] & c
  \end{tikzcd}
\]
In any case, keep in mind that a functor is what running under a diagram.

\begin{example}[Categorial axioms]
  Let us express the categorial axioms in a diagrammatic vest. Let
  \(\mathcal{C}\) be a category and \(x\) an object of \(\mathcal{C}\). The fact
  that \(\id_x\) the identity of \(x\) can be translated as follows:
  the diagrams
  \begin{equation}\begin{tikzcd}[row sep=small]
      & x \ar["{\id_x}", dd] \\
      a \ar["f", ur] \ar["f", dr, swap]  \\
      & x &
    \end{tikzcd} \quad
    \begin{tikzcd}[row sep=small]
      x \ar["{\id_x}", dd, swap] \ar["g", dr] & \\
      &  b \\
      x \ar["g", ur, swap] &
    \end{tikzcd}
    \label{diag:IdPropDiag}\end{equation}
  commute for every \(a\) and \(b\) objects and \(f\) and \(g\)
  morphisms of \(\mathcal{C}\). Associativity can be rephrased by saying:
  \[
    \begin{tikzcd}
      a \ar["f", r] \ar["{gf}", swap, dr] & b \ar["{hg}", dr] \\
      & c \ar["h", swap, r] & d
    \end{tikzcd}
  \]
  commutes for every \(a\), \(b\), \(c\) and \(d\) objects and \(f\),
  \(g\) and \(h\) morphisms in \(\mathcal{C}\).
\end{example}
% Cause diagrams are supposed to be drawn, we shall assume some
% conventions. The spatial placement of nodes and the shape of arrows
% is a matter of aesthetics, just pursue clarity. Do not burden your
% diagrams: though arrows can be composed, we usually do not draw
% them; similarly, neither identities shall be drawn.

\begin{example}[Homotopy]
  Consider two continuous functions \(f, g : X \to Y\); so we are in
  \(\Top\). We say \(f\) is {\em homotopic} to \(g\) whenever there is
  some continuous function \(H : X \times I \to Y\), the {\em
    homotopy}, such that
  \begin{align*}
    & H(x, 0) = f(x) \text{ for every } x \in X \\
    & H(x, 1) = g(x) \text{ for every } x \in X
  \end{align*}
  How do we write this in diagrammatic fashion? We can consider two
  devices: \(i_0 : X \to X \times I\) defined by \(i_0(x) = (x, 0)\)
  and \(i_1 : X \to X \times I\) defined by \(i_1(x) = (x, 1)\). In
  fact, the two conditions above can be equivalently stated as: the
  diagram
  \[
    \begin{tikzcd}
      X \ar["{i_0}", r] \ar["f", dr, swap] & X \times I \ar["H", d] &
      X \ar["{i_1}", l,
      swap] \ar["g", dl] \\
      & Y
    \end{tikzcd}
  \]
  commutes.
\end{example}


\begin{example}[Semigroup axioms]
  A {\em semigroup} is a set \(X\) together with a function
  \(\mu : X \times X \to X\) which is associative, that is
  \[
    \mu(\mu(a, b), c) = \mu(a, \mu(b, c)) \ \text{for every } a, b, c
    \in X .
  \]
  The aim of this example is to see how can we put in diagrams all
  this. We have a triple of elements, to start with,
  \((a, b, c) \in X \times X \times X\). On the left side of the
  equality above, \(a\) and \(b\) are multiplied first, and the result
  is multiplied with \(c\):
  \[
    \begin{tikzcd}[row sep=tiny, column sep=small]
      X \times X \times X \ar[r] & X \times X \ar[r] & X \\
      (a, b, c) \ar[r] & (\mu(a, b), c) \ar[r] & \mu(\mu(a, b), c)
    \end{tikzcd}
  \]
  It is best we some effort in naming these functions. While it is
  clear that \(X \times X \to X\) is our \(\mu\), how do we write
  \(X \times X \times X \to X\)? There is notation for it:
  \(\mu \times \id_X\).\footnote{In general, if you have two functions
    \(f : A_1 \to A_2\) and \(g : B_1 \to B_2\), the function
    \(f \times g : A_1 \times B_1 \to A_2 \times B_2\) is the one
    defined by \(f \times g (a, b) = (f(a), g(b))\).} Instead, on the
  other side, \(b\) and \(c\) are multiplied first, and then \(a\) is
  multiplied to their product:
  \[
    \begin{tikzcd}[row sep=tiny, column sep=small]
      X \times X \times X \ar[r] & X \times X \ar[r] & X \\
      (a, b, c) \ar[r] & (a, \mu(b, c)) \ar[r] & \mu(a, \mu(b, c))
    \end{tikzcd}
  \]
  The first function is \(\id_X \times \mu\) and the second one is
  \(\mu\). Thus the equality of the definition of semigroup is
  equivalent to the fact that the diagram
  \[
    \begin{tikzcd}
      X \times X \times X \ar["{\mu \times \id_X}", r] \ar["{\id_X \times \mu}", d, swap] & X \times X \ar["\mu", d] \\
      X \times X \ar["\mu", r, swap] & X
    \end{tikzcd}
  \]
  commutes.
\end{example}

\begin{exercise}
  Recall that a monoid is a semigroup \((X, \mu)\) with \(e \in X\)
  such that \(\mu(x, e) = \mu(e, x) = x\) for every \(x \in X\). We
  usually write a monoid as a triple \((X, \mu, e)\). A {\em group} is
  a monoid \((G, \mu, e)\) such that for every \(x \in G\) there
  exists \(y \in G\) such that \(\mu(x, y) = \mu(y, x) = e\). Rewrite
  these structures using commutative diagrams. (The work about
  associativity is already done, so you should focus how to express
  the property of the identity in a monoid and the property of
  inversion for groups.)  Also recall that a {\em monoid
    homomorphism}, then, from a monoid \((X, \mu, e_X)\) to a monoid
  \((Y, \lambda, e_Y)\) is any function \(f : X \to Y\) such that
  \(f(\mu(a, b)) = \lambda (f(a), f(b))\) for every \(a, b \in X\) and
  \(f(e_X) = e_Y\). Use commutative diagrams. Observe that group
  homomorphisms are defined by requiring to preserve multiplication,
  whereas the the preservation of identities can be deduced.
\end{exercise}



\section{Duality}

\begin{construction}
For \(\mathcal{C}\) a category, its {\em dual} (or {\em opposite}) category
is denoted \(\opcat C\) and is described as follows:
\begin{itemize}
\item The objects are the same of \(\mathcal{C}\), that is
  \(\obj{\opcat C} = \obj{\mathcal{C}}\).
\item The morphism \(a \to b\) in \(\opcat C\) are the morphisms
  \(b \to a\) in \(\mathcal{C}\). In other words,
  \(\opcat C (a, b) = \mathcal{C}(b, a)\).
\item For
  \(\begin{tikzcd}[column sep=small, cramped] a \ar["{f}", r] & b
    \ar["{g}", r] & c \end{tikzcd}\) in \(\opcat C\), i.e.
  \(\begin{tikzcd}[column sep=small, cramped] c \ar["{g}", r] & b
    \ar["{f}", r] & a \end{tikzcd}\) in \(\mathcal{C}\), the composition is
  given by
  \[
    gf := fg .
  \]
\end{itemize}
The dual category has the same objects of the original one. They have
the same morphisms too, but have their verses reversed: a morphism
\(a \to b\) of \(\opcat C\) can be found among morphisms \(b \to a\) in
\(\mathcal{C}\). This has consequence on the composition as well: the
composition in \(\opcat C\) is that of \(\mathcal{C}\) but with arguments
reversed. The proof of the categorial properties for dual categories is
straightforward as it reduces to computation in the original one.
\end{construction}

\begin{example}[Dual prosets]
  If \((A, \le)\) is a proset, then (categorially speaking) its dual is
  the proset \((A, \ge)\) where for \(a, b \in A\) we define \(b \ge a\)
  to be \(a \le b\).
\end{example}

% \begin{exercise}
%   Consider a single object category \(\mathcal{G}\), that is a monoid. What
%   is \(\opcat G\)? What is \(\opcat G\) if \(\mathcal{G}\) is a group?
% \end{exercise}

\begin{example}
  Let \(\mathcal{C}\) be a monoid, that is a single object category, and look
  at \(\opcat C\). The unique part to make explicit is the
  composition. Given \(m\) and \(n\) in \(\opcat C\)
  \[
    \begin{tikzcd}[column sep=small]
      \bullet \ar["{m}", r] & \bullet \ar["{n}", r] & \bullet
    \end{tikzcd}
  \]
  then the composite \(n m\) in \(\opcat C\) is defined to be \(m n\).
\end{example}

It may seem hard to believe, but duality is one of the biggest conquest
of Category Theory. To understand this, we provide examples.

\begin{example}
  The set of natural numbers \(\naturals\) has the order relation of
  divisibility, that we denote \(\divides\): this poset is a category,
  as we have seen. From Group Theory, we know that for every
  \(m, n \in \naturals\) such that \(m \divides n\) there is a
  homomorphism
  \begin{align*}
    & f_{m, n} : \integers_n \to \integers_m \\
    & f_{m, n} \left([a]_n\right) := [a]_m .
  \end{align*}
  Our hopes to carve out a functor \((\naturals, \divides) \to \Grp\) are
  quickly shuttered: from \(m \divides n\) we have a homomorphism
  \(\integers_n \to \integers_m\), not \(\integers_m \to \integers_n\)
  (observe the order of \(m\) and \(n\)!). Instead, this example offers
  us a nice functor:
  \[
    F : \op{(\naturals, \divides)} \to \Grp
  \]
  that maps naturals \(n\) to groups \(\integers_n\) and
  \(m \op\divides n\), that is \(n \divides m\), to the group
  homomorphism \(f_{m,n}\) defined above.
\end{example}

Traditionally, functors of Definition~\ref{definition:functors} are
called \q{covariant}, because there are {\em contra}variant functors
too. For if \(\mathcal{C}\) and \(\mathcal{D}\) are categories, a {\em
  contravariant functor} from \(\mathcal{C}\) to \(\mathcal{D}\) is just a functor
\(\op{\mathcal{C}} \to \mathcal{D}\).

Let us unpack a functor \(F : \op{\mathcal{C}} \to \mathcal{D}\). They map objects
to objects and morphisms \(f : a \to b\) of \(\op{\mathcal{C}}\) to morphisms
\(F(f) : F(a) \to F(b)\) of \(\mathcal{D}\). But, remembering how dual
categories are defined, what \(F\) actually does is this:
\begin{quotation}
  it maps objects of \(\mathcal{C}\) to objects of \(\mathcal{D}\) and morphisms
  \(f : b \to a\) of \(\mathcal{C}\) to morphisms \(F(f) : F(a) \to F(b)\) of
  \(\mathcal{D}\) (mind that \(a\) and \(b\) have their roles flipped).
\end{quotation}
Now, what about functoriality axioms? Neither with identities \(F\) does
something different and the composite \(gf\) of \(\op{\mathcal{C}}\) is
mapped to the composite \(F(g)F(f)\) of \(\mathcal{D}\). Again by definition
of dual categories, this can be translated as follows:
\begin{quotation}
  the composite \(fg\) of \(\mathcal{C}\) is mapped to \(F(g)F(f)\) (notice
  here how \(f\) and \(g\) have their places switched).
\end{quotation}

% What does duality mean for diagrammatic reasoning?  Commuting triangles
% \[\begin{tikzcd}[row sep=small]
%     & b \ar["g", dd] \\
%     a \ar["f", ur] \ar["h", dr, swap] \\
%     & c
%   \end{tikzcd}\] of \(\mathcal{C}\) are exactly commuting triangles
% \[\begin{tikzcd}[row sep=small]
%     & b \ar["f", dl, swap] \\
%     a \\
%     & c \ar["h", ul] \ar["g", uu, swap]
%   \end{tikzcd}\] in \(\opcat C\).

% The concept of duality for categories has one important consequence on
% statements written in a \q{categorial language}. We do not need to be
% fully precise here: they are statements written in a sensible way using
% the usual logical connectives, names for objects, names for morphisms
% and quantifiers acting on such names.

% \begin{example}
%   If we have a morphism \(f : a \to b\) in some category \(\mathcal{C}\),
%   consider the statement
%   \begin{quotation}
%     For every {\color{red!65!black} object \(c\) of \(\mathcal{C}\)} and
%     {\color{blue!65!black} morphisms \(g_1, g_2 : c \to a\) in
%       \(\mathcal{C}\)}, if {\color{green!65!black} \(f g_1 = f g_2\)} then
%     {\color{orange!65!black}\(g_1 = g_2\)}.
%   \end{quotation}
%   If you remember, it is just said that \(f\) is a monomorphism. We
%   operate a translation that doesn't modify the truth of the sentence:
%   that is, if it is true, it remains so; it is false, it remains false.
%   \begin{quotation}
%     For every {\color{red!65!black} object \(c\) of \(\opcat C\)} and
%     {\color{blue!65!black} morphisms \(g_1, g_2 : a \to c\) in
%       \(\opcat C\)}, if {\color{green!65!black} \(g_1 f = g_2 f\)} then
%     {\color{orange!65!black}\(g_1 = g_2\)}.
%   \end{quotation}
%   If we regard \(f\) as a morphism \(b \to a\) of \(\opcat C\), then
%   \(f\) is an epimorphism in \(\opcat C\).
% \end{example}

% Let us try to settle this explicitly: if we have a categorial statement
% \(p\), the dual of \(p\) --- we may call \(\op{p}\) --- is the statement
% obtained from \(p\) keeping the connectives and the quantifiers of
% \(p\), whereas the other parts are replaced by their dual counterparts.

% \begin{example}
%   \NotaInterna{Anticipate products and co-products\dots{}}
% \end{example}



\section{Bifunctors}

\begin{construction}
  Given two categories \(\mathcal{C}\) and \(\mathcal{D}\), the {\em product
    category} \(\mathcal{C} \times \mathcal{D}\) is the category where:
  \begin{itemize}[leftmargin=*]
  \item The objects are exactly pairs
    \((x, y) \in \obj{\mathcal{C}} \times \obj{\mathcal{D}}\).
  \item For every \((x, x'), (y, y') \in \obj{\mathcal{C} \times \mathcal{D}}\), a
    morphism from \((x, x')\) to \((y, y')\) is defined to be a pair
    \[
      \left(
        \begin{tikzcd}[row sep=small, cramped] x \ar["{f}", d, swap] \\
          x' \end{tikzcd} ,
        \begin{tikzcd}[row sep=small, cramped] y \ar["{g}", d] \\
          y' \end{tikzcd} \right)
    \]
  \item The composition of morphisms
    \[
      \begin{tikzcd}
        (x, x') \ar["{(f, f')}", r] & (y, y') \ar["{(g, g')}", r] & (z,
        z')
      \end{tikzcd}
    \]
    is defined to be \((gf, g'f')\).
  \end{itemize}
\end{construction}

\begin{example}
  \NotaInterna{Functoriality of the product of sets.}
\end{example}

In a locally small \NotaInterna{have we defined somewhere that?}
category \(\mathcal{C}\), take two morphisms
\[
  \begin{tikzcd}
    a_1 \\
    b_1 \ar["g_1", u]
  \end{tikzcd} \quad\text{and}\quad \begin{tikzcd}
    a_2 \ar["g_2", d] \\
    b_2
  \end{tikzcd}
\]
For \(f : a_1 \to a_2\) we have the composite
\(g_2 f g_1 : b_1 \to b_2\), that is a function
\[\mathcal{C} (a_1, a_2) \to \mathcal{C} (b_1, b_2) .\]
We will refer to this function using lambda calculus notation:
\(\lambda f . g_2 f g_1\). If we have
\[
  \begin{tikzcd}[row sep=small]
    a_1 \ar["f", r] & a_2 \ar["g_2", d] \\
    b_1 \ar["g_1", u] & b_2 \ar["h_2", d] \\
    c_1 \ar["h_1", u] & c_2
  \end{tikzcd}
\]
we can say that
\[
  \lambda f . (h_2 g_2) f (g_1 h_1) = (\lambda f' . h_2 f' h_1) (\lambda
  f . g_2 f g_1) ,\] which can be derived by uniquely using the
associativity of the composition. Another remarkable property can be
obtained when \(a_1 = b_1\), \(g_1 = \id_{a_1}\), \(a_2 = b_2\) and
\(g_2 = \id_{a_2}\):
\[
  \lambda f . \id_{a_2} f \id_{a_1} = \lambda f . f = \id_{\mathcal{C} (a_1,
    b_1)} .
\]
There is functoriality, to understand that, we need to package all this
machinery in one functor.

\begin{definition}
  If \(\mathcal{C}\) is a locally small category, the functor
  \[
    \hom_{\mathcal{C}} : \opcat C \times \mathcal{C} \to \Set
  \]
  takes every \((x, y) \in \obj{\opcat C \times \mathcal{C}}\) to
  \(\mathcal{C} (x, y)\) and on morphisms in \(\opcat C \times \mathcal{C}\) is
  given as follows:
  \[
    \hom_{\mathcal{C}} \left(\begin{tikzcd}[row sep=small] a_1 \\ b_1
        \ar["g_1", u] \end{tikzcd}, \begin{tikzcd}[row sep=small] a_2
        \ar["g_2", d] \\ b_2 \end{tikzcd}\right)
    \left(\begin{tikzcd}[column sep=small] a_1 \ar["f", r] &
        a_2 \end{tikzcd}\right) := g_2 f g_1 .
  \]
\end{definition}

Assume we have a functor \(F : \mathcal{C} \times \mathcal{D} \to \mathcal{E}\). In
total analogy to what happens with functions in two variables, we can
curry functors as well. Fixed some \(a \in \obj{\mathcal{C}}\), we introduce
the functor
\[
  F(a, -) : \mathcal{D} \to \mathcal{E}
\]
which takes \(b \in \obj{\mathcal{D}}\) to \(F(a, b)\) and \(g : b \to b'\)
to \(F(\id_a, g) : F(a, b) \to F(a, b')\). Similarly, assigned one
\(b \in \obj{\mathcal{D}}\), we introduce the functor
\[
  F(-, b) : \mathcal{C} \to \mathcal{E}
\]
which takes \(a \in \obj{\mathcal{C}}\) to \(F(a, b)\) and \(f : a \to a'\)
to \(F(f, \id_b) : F(a, b) \to F(a', b)\).

\begin{definition}\label{definition:hom-functors}
  Let \(\mathcal{C}\) be a locally small category. From
  \(\hom_{\mathcal{C}} : \opcat C \times \mathcal{C} \to \Set\) we have the
  functors
  \begin{align*}
    & \mathcal{C}(a, -) := \hom_{\mathcal{C}}(a, -) : \mathcal{C} \to \Set \\
    & \mathcal{C}(-, b) := \hom_{\mathcal{C}}(-, b) : \opcat C \to \Set
  \end{align*}
\end{definition}

We will meet again these functors later; for the moment consider some
simple example.

\begin{example}
  If \(\mathcal{G}\) is a monoid, a category whose unique object is written
  as \(\bullet\), then the functor
  \(\mathcal{G}(\bullet, -) : \mathcal{G} \to \Set\) is an action on the set
  \(M := \mathcal{G}(\bullet, \bullet)\). Unpacking a bit, it maps a morphism
  \(m : \bullet \to \bullet\) to the function
  \begin{align*}
    & M \to M \\
    & g \mapsto mg
  \end{align*}
  In other words, \(\mathcal{G} (\bullet, -)\) is the action of the monoid
  \(M\) on the set \(M\) by multiplication on the left. It is not
  difficult at this point to understand what
  \(\mathcal{G}(-, \bullet) : \opcat C \to \Set\) is.
\end{example}

In some texts you might find written \(f_\ast\) for
\(\mathcal{C}(a, f) = \mathcal{C}(\id_a, f)\) and \(g^\ast\) instead of
\(\mathcal{C}(g, b) = \mathcal{C}(g, \id_b)\).

\begin{example}[The dual space]\label{example:DualSpace}
  If \(V\) is vector space on a field \(k\), we write \(V^\ast\) for its
  dual, that is the set of linear maps \(V \to k\). (As a matter of
  fact, this is a vector space on \(k\) too, but let us forget that for
  the moment.) If \(f : V \to W\) is a linear map, then we have the
  function
  \[
    f^\ast : W^\ast \to V^\ast
  \]
  sending linear maps \(g : W \to k\) to linear maps \(gf : V \to
  k\). (Well, also \(f^\ast\) is linear, but let us ignore this.) We
  have a functor \(\op\Vect_k \to \Set\) which is exactly
  \(\Vect_k(-, k)\)!
\end{example}



\section{Natural transformations}

Given two functors \(F, G : \mathcal{C} \to \mathcal{D}\), a {\em transformation}
from \(F\) to \(G\) is a family of morphisms of \(\mathcal{D}\)
\[
  \begin{tikzcd}[column sep=small]
    F(x) \ar["{\eta_x}", r] & G(x)
  \end{tikzcd}
  \text{ for } x \in \obj{\mathcal{C}}
\]
The idea behind is measuring the difference of two parallel functor by
the unique means we have, viz morphisms.

In general, we stick to the following convention: if \(\eta\) is the
name of a transformation from \(F\) to \(G\), then \(\eta_x\) indicates
the component \(F(x) \to G(x)\) of the transformation.

We are not interested in all transformations, of course.

\begin{definition}[Natural transformations]
  A transformation \(\eta\) from a functor \(F : \mathcal{C} \to \mathcal{D}\) to
  a functor \(G : \mathcal{C} \to \mathcal{D}\) is said to be {\em natural}
  whenever for every \(a, b \in \obj{\mathcal{C}}\) and
  \(f \in \mathcal{C}(a, b)\) the square
  \[
    \begin{tikzcd}
      F(a) \ar["{\eta_a}", r] \ar["{F(f)}", d, swap] & G(a) \ar["{G(f)}", d] \\
      F(b) \ar["{\eta_b}", r, swap] & G(b)
    \end{tikzcd}
  \]
  commutes. This property is called \q{naturality} of \(\eta\).
\end{definition}

We write \(\eta : F \tto G\) to say \(\eta\) is a natural transformation
from \(F\) to \(G\). If we want to -- or need to -- give the picture of
the whole context we might go for something like this:
\[
  \naturaltr{\eta}{\mathcal{C}}{\mathcal{D}}{F}{G}
\]

Sometimes, we introduce a family of morphisms \(\eta_x : Fx \to Gx\)
parameterised by the objects of the domain of \(F\) and \(G\) and say
such morphisms are {\em natural in \(x\)}. This jargon motivated by
practice: one isolates the single maps \(Fx \to Gx\) and afterwards learns
they are natural.

\begin{example}
  Let \(V\) be a vector space on some field \(k\). Recall that
  \(V^\ast\) is the vector space on \(k\) of the linear maps
  \(f : V \to k\). Introduce now the linear map
  \[
    e_V : V \to V^{\ast\ast}
  \]
  taking \(\mathbf v \in V\) to the linear map \(V^\ast \to k\),
  \(f \mapsto f(\mathbf v)\). If we want to see if \(e_V\) is natural in
  \(V\), we need to make sure categories and functors are
  involved. Well, it seems the identity functor \(\id_{\Vect_k}\) is
  involved, hence we want another functor \(\Vect_k \to \Vect_k\). In fact
  there is a functor
  \[
    (\phantom\square)^{\ast\ast} : \Vect_k \to \Vect_k
  \]
  which takes \(V\) to \(V^{\ast\ast}\) and sends \(f : V \to W\) to
  \begin{align*}
    & f^{\ast\ast} : V^{\ast\ast} \to W^{\ast\ast} \\
    & f^{\ast\ast} (g) = g f^\ast
  \end{align*}
  (Compare with Example~\ref{example:DualSpace}.)
\end{example}

Again Linear Algebra.

\begin{example}[Determinant of matrices]
  Fix one \(n \in \naturals\) and a commutative ring (with unity)
  \(R\). Then the matrices of type \(n \times n\) with entries in \(R\)
  together with the usual matrix multiplication form a monoid. We write
  such object as \(M_n(R)\). The {\em determinant} is a monoid
  homomorphism
  \[
    \det_{n, R} : M_n(R) \to R .
  \]
  To complete the picture, we put categories and functors. One functor
  is the forgetful functor \(U : \CRing \to \Mon\) and the other is
  \(M_n : \CRing \to \Mon\), which takes
  \begin{itemize}
  \item commutative rings \(R\) to \(M_n(R)\)
  \item ring homomorphism \(f : R \to S\) to the monoid homomorphism
    \(M_n(f)\) sending a matrix \(A\) to the matrix with \(f(A_{ij})\)
    as entry at position \((i, j)\).
  \end{itemize}
  That said, we can easily verify \(\det_{n, R}\) is natural in \(R\).
\end{example}

As you may have already surmised, natural transformation are morphisms
between functors. Indeed there is a nation of compositionality
satisfying categorial axioms.

Consider two consecutive natural transformations
\[
  \begin{tikzcd}
    \mathcal{C} \ar["F"{name=F}, rr, bend left=70] \ar["G"{name=G,
      description}, rr] \ar["H"{name=H}, rr, bend right=70, swap] & &
    \mathcal{D} \ar["\eta", natural, from=F, to=G] \ar["\theta", natural, from=G,
    to=H]
  \end{tikzcd}
\]
We can define \(\theta \eta\) to be the family of morphisms
\[
  (\theta \eta)_x := \theta_x \eta_x : Fx \to Hx \quad\text{for } x \in \obj{\mathcal{C}} .
\]
It is simple to verify that \((\theta \eta)_x\) is natural in \(x\) by simply
looking at the diagram
\[
  \begin{tikzcd}
    Fa \ar["{Ff}", d, swap] \ar["{\eta_a}", r] & Ga \ar["{Gf}", d]
    \ar["{\theta_a}", r] & Ha \ar["{Hf}", d] \\
    Fb \ar["{\eta_b}", r, swap] & Gb \ar["{\theta_b}", r, swap] & Hb
  \end{tikzcd}
\]
Such composition is associative; moreover, for every functor
\(F: \mathcal{C} \to \mathcal{D}\) there is the natural transformation
\(\id_F : F \tto F\) with components \(\id_{Fx} : Fx \to Fx\), for
\(x \in \obj{\mathcal{C}}\). Such natural transformations are identities in
categorial sense.

\begin{definition}
  Given categories \(F, G : \mathcal{C} \to \mathcal{D}\), we write
  \([\mathcal{C}, \mathcal{D}]\) the {\em functor category}, that is the category
  of functors and natural transformations depicted earlier. In
  particular, a {\em natural isomorphism} is an isomorphism in
  \([\mathcal{C}, \mathcal{D}]\).
\end{definition}

% All this suggests to, given two categories \(\mathcal{C}\) and \(\mathcal{D}\),
% form a category with functors \(\mathcal{C} \to \mathcal{D}\) as objects and natural
% transformations as morphism, them being composable as explained
% above. [\dots{}]

% \NotaInterna{Consider \url{https://mathoverflow.net/q/39073}\dots{}}

There is a simple characterisation of isomorphisms of functor categories.

\begin{proposition}
  A natural transformation is a natural isomorphism if and only if its
  components are all isomorphisms.
\end{proposition}

\begin{exercise}
  Prove the previous Proposition.
\end{exercise}

Sometimes we will need to work with natural transformations between
bifunctors. Verifying naturality in two variables can be difficult
sometimes though. Luckily you can split the work in two: check
naturality in the variables separately.

\begin{proposition}\label{proposition:NaturalityInToVariables}
  Let \(F, G : \mathcal{C}_1 \times \mathcal{C}_2 \to \mathcal{D}\) be functors and a transformation \(\eta : F
  \Rightarrow G\) a transformation. Then the following are equivalent:
  \begin{enumerate}[leftmargin=*]
  \item \(\eta\) is natural.
  \item For every \(a \in \obj{\mathcal{C}_1}\), the morphism
    \(\eta_{a, b} : F(a, b) \to G(a, b)\) is natural in \(b\) and for every
    \(b \in \obj{\mathcal{C}_2}\), the morphism
    \(\eta_{a, b} : F(a, b) \to G(a, b)\) is natural in \(a\).
  \end{enumerate}
\end{proposition}

\begin{proof}
  
\end{proof}

% \begin{example}
%   Consider the bifunctor
%   \[
%     (\times) : \Set \times \Set \to \Set
%   \]
%   that sends pairs \((A, B)\) to cartesian products \(A \times B\) and pairs
%   \((f : A \to A', g : B \to B')\) to
%   \begin{align*}
%     & f \times g : A \times B \to A' \times B' \\
%     & (f \times g)(a, b) := (f(a), g(b))
%   \end{align*}
%   Now consider the bifunctor
%   \[
%     (\tilde\times) : \Set \times \Set \to \Set
%   \]
%   defined which is \(\times\) but swaps arguments: sends \((A, B)\) to
%   \(B \times A\) and \((f, g)\) to \(g \times f\). Prove now that
%   \begin{align*}
%     & \swap_{A, B} : A \times B \to B \times A \\
%     & \swap_{A, B} (a, b) := (b, a)
%   \end{align*}
%   is natural in \(A\) and \(B\), that is \(\swap : (\times) \Rightarrow (\tilde\times)\). 
% \end{example}

% \begin{example}
%   Consider for \(A, B, C \in \obj\Set\) the function
%   \begin{align*}
%     & (\hole)^\sharp : \Set(A \times B, C) \to \Set(A, \Set(B, C)) \\
%     & f^\sharp := \lambda a . \lambda b . f(a, b)
%   \end{align*}
%   and verify that it is natural in \(A, B, C \in \obj\Set\). To do so,
%   make sure you get the whole picture. One functor is
%   \[
%     F : \op\Set \times \op\Set \times \Set \to \Set
%   \]
%   which takes \((A, B, C)\) to \(\Set (A \times B, C)\) and \((f : A' \to A, g
%   : B' \to B, h : C \to C')\) to the function
%   \begin{align*}
%     & F(f, g, h) : \Set(A \times B, C) \to \Set( A' \times B', C') \\
%     & F(f, g, h)(k) := hk(f \times g)
%   \end{align*}
%   The other functor is
%   \[
%     G : \op\Set \times \op\Set \times \Set \to \Set
%   \]
%   which takes \((A, B, C)\) to \(\Set (A, \Set(B, C))\) and \((f : A' \to A, g
%   : B' \to B, h : C \to C')\) to the function
%   \begin{align*}
%     & G(f, g, h) : \Set(A, \Set(B, C)) \to \Set(A', \Set(B', C')) \\
%     & G(f, g, h)(k) := \lambda a' . \lambda b' . hk(f a')(g b')
%   \end{align*}
% \end{example}


\section{Equivalence of categories}

It might seem the discussion about ``sameness'' categories can be easily
liquidated: two categories \(\mathcal{C}\) and \(\mathcal{D}\) are the same whenever they are isomorphic, in
the sense there are functors \(F : \mathcal{C} \to \mathcal{D}\) and \(G : \mathcal{D} \to
C\) such that \(GF = \id_{\mathcal{C}}\) and \(FG = \id_{\mathcal{D}}\). Well,
it's not that easy.

\begin{definition}
  We say a category \(\mathcal{C}\) is {\em equivalent} to a category
  \(\mathcal{D}\), and write \(\mathcal{C} \simeq \mathcal{D}\), when there are functors
  \[
    \begin{tikzcd}[column sep=small, cramped] \mathcal{C} \ar["{F}", r, shift
      left] & \mathcal{D} \ar["{G}", l, shift left] \end{tikzcd}
  \]
  and two natural isomorphisms
  \[
    \eta : \id_{\mathcal{C}} \Rightarrow GF \quad\text{and}\quad \epsilon : FG \Rightarrow \id_{\mathcal{D}} .
  \]
  The tuple \((F, G, \eta, \epsilon)\) is called {\em equivalence}. In other
  books, you might find written that an equivalence is just a functor
  \(F : \mathcal{C} \to \mathcal{D}\) for which there exist
  \(G : \mathcal{D} \to \mathcal{C}\) and natural isomorphisms
  \(\eta : \id_{\mathcal{C}} \Rightarrow GF\) and \(\epsilon : FG \Rightarrow \id_{\mathcal{D}}\).
\end{definition}

Since \(\eta\) and \(\epsilon\) are natural isomorphisms, then also
\((G, F, \inv\epsilon, \inv\eta)\) is an equivalence. However, observe that the
equivalence is \(\mathcal{D} \simeq \mathcal{C}\) and we have never written that the
equivalence of categories is an equivalence relation. Yes, it is, but
that comes after some work.

Here is a reformulation of the concept of equivalence of categories.

\begin{proposition}
  Let \((F, G, \eta, \epsilon)\) be an equivalence between \(\mathcal{C}\) and
  \(\mathcal{D}\). Then \(F\) are \(G\) are full, faithful and essentially
  surjective on objects. Conversely -- if you accept the {\sc Axiom of
    Choice} -- a full, faithful and essentially surjective functor
  \(F : \mathcal{C} \to \mathcal{D}\) gives an equivalence
  \((F, G, \eta, \epsilon)\) of categories such that
  \begin{align}
    & (\epsilon F)(F \eta) = \id_F \label{eqn:TriId1} \\
    & (G \epsilon) (\eta G) = \id_G \label{eqn:TriId2}
  \end{align}
\end{proposition}

\NotaInterna{Talk about the natural transformations
  in~\eqref{eqn:TriId1} and~\eqref{eqn:TriId2}.}

This is the reason why sometimes we refer to full, faithfull and
essentially surjective functors as equivalences.

\begin{proof}
  (\(\Rightarrow\)) We prove only that \(F\) is full, faithful and essentially
  surjective on objects, since for \(G\) the argument is similar. Let
  \(f, f' \in \mathcal{C} (a, b)\) such that \(Ff = Ff'\). By naturality of
  \(\eta\), the squares in
  \[
    \begin{tikzcd}
      a \ar["{f}", d, swap, shift right] \ar["{f'}", d, shift left] \ar["{\eta_a}", r] & GFa \ar["{GFf = GFf'}", d] \\
      b \ar["{\eta_b}", r, swap] & GFb
    \end{tikzcd}
  \]
  commutes. The equality \(f = f'\) follows immediately because
  \[
    f = \inv \eta_b (GF f) \eta_a = \inv\eta_b (GF f') \eta_a = f' .
  \]
  Now let \(g \in \mathcal{D}(Fa, Fb)\) and pick
  \(f := \inv \eta_b(Gg) \eta_a \in \mathcal{C}(a, b)\). Again by naturality of
  \(\eta\), we have the following commutative square
  \[
    \begin{tikzcd}
      a \ar["{f}", d, swap] \ar["{\eta_a}", r] & GFa \ar["{GFf}", d] \\
      b \ar["{\eta_b}", r, swap] & GFb
    \end{tikzcd}
  \]
  In this case \(G(Ff) = \eta_b f \inv\eta_a = Gg\), which implies
  \(Ff = g\) because \(G\) is faithful. To conclude, observe also that
  for \(d \in \obj{\mathcal{D}}\), we have \(F(Gd) \cong d\) thanks to \(\epsilon\).

  (\(\Leftarrow\)) Assume you have a functor \(F : \mathcal{C} \to \mathcal{D}\) full,
  faithful and essentially surjective: from this we will construct
  \(G : \mathcal{D} \to \mathcal{C}\) and natural isomorphisms
  \(\eta : \id_{\mathcal{C}} \Rightarrow GF\) and \(\epsilon : FG \Rightarrow \id_{\mathcal{D}}\).
  \begin{itemize}
  \item Since \(F\) is essentially surjective, for every
    \(d \in \obj{\mathcal{D}}\) choose one object in \(\mathcal{C}\) that we name
    \(Gd\) and such that \(F(G d) \cong d\). Since we need this notation
    later, write \(\epsilon_d : F(G d) \to d\) one of the morphisms realising the
    isomorphism \(F(G d) \cong d\). Now for every \(g : a \to b\) in
    \(\mathcal{D}\), let \(Gg : Ga \to Gb\) the morphism such that
    \[
      F(Gg) := \inv\epsilon_b g \epsilon_a
    \]
    \[
      \begin{tikzcd}
        a \ar["{g}", d, swap] & FGa \ar["{\epsilon_a}", l, swap] \ar["{FGg}", d] \\
        b \ar["{\inv\epsilon_b}", r, swap] & FGb
      \end{tikzcd}
    \]
    We shall now verify \(G\) gives arise to a functor
    \(\mathcal{D} \to \mathcal{C}\). First of all,
    \[
      F\left( G\id_a \right) = \inv\epsilon_a \id_a \epsilon_a =
      F\left(\id_{Ga}\right) .
    \]
    Being \(F\) faithful, we can derive \(G \id_a = \id_{Ga}\). Now take
    two composable morphisms
    \(\begin{tikzcd}[column sep=small, cramped] a \ar["{f}", r] & b
      \ar["{g}", r] & c \end{tikzcd}\).
    \begin{align*}
      F(G(gf)) &= \inv\epsilon_c (gf) \epsilon_a = \left( \inv\epsilon_c g \epsilon_b \right) \left(
                 \inv\epsilon_b f \epsilon_a \right) = \\
               &= F(Gg) F(Gf) = F(Gg Gf) .
    \end{align*}
    Again by faithfulness of \(F\), conclude that \(G(gf) = Gg Gf\).
  \item We have already introduced a family of isomorphisms
    \(\epsilon_a : FGa \to a\). These morphisms are natural in
    \(a \in \obj{\mathcal{D}}\) just by construction of \(G\) on morphisms.
  \item Now, for every \(x \in \obj{\mathcal{C}}\) we shall introduce
    \(\eta_x : x \to GFx\) which is natural in \(x\). Take the morphism
    \(\begin{tikzcd}[column sep=small, cramped] Fx
      \ar["{\inv\epsilon_{Fx}}", r] & FGFx \end{tikzcd}\) in \(\mathcal{D}\):
    being \(F\) full and faithful, there is a unique
    \(\eta_x : x \to GFx\) such that \(F\eta_x =
    \inv\epsilon_{Fx}\). Observe that~\eqref{eqn:TriId1} is
    proved. Observe also that \(\eta_x\) is an isomorphism, since \(F\)
    is full and faithful. Consider the square
    \[
      \begin{tikzcd}
        x \ar["{f}", d, swap] \ar["{\eta_x}", r] & GFx \ar["{GFf}", d] \\
        y \ar["{\eta_y}", r, swap] & GFy
      \end{tikzcd}
    \]
    We have
    \begin{align*}
      & F\left((GFf) \eta_x\right) = (FGFf) (F\eta_x) = (FGFf) \inv\epsilon_{Fx} \\
      & F\left(\eta_yf\right) = (F\eta_y) (Ff) = \inv \epsilon_{Fy} (Ff)
    \end{align*}
    The naturality of \(\eta\) derives from the naturality of \(\epsilon\)
    \[
      \begin{tikzcd}
        FGFx \ar["{FGFf}", d, swap] \ar["{\epsilon_{Fx}}", r] & Fx \ar["{Ff}", d] \\
        FGFy \ar["{\epsilon_{Fy}}", r, swap] & Fy
      \end{tikzcd}
    \]
    and from the fact that \(F\) is faithful.
    
  \item By how the functor \(G\) is introduced, we have
    \[
      F\left(G\epsilon_a\right) = \inv\epsilon_a \epsilon_a \epsilon_{FGa} = \epsilon_{FGa}
    \]
    On the other hand, by how \(\eta\) is introduced, we have
    \(F(\eta_{Ga}) = \inv\epsilon_{FGa}\), from which we deduce that
    \[
      F\left( \inv\eta_{Ga} \right) = \epsilon_{FGa}.
    \]
    To derive~\eqref{eqn:TriId2} use the fact that \(F\) is
    faithful.\qedhere
  \end{itemize}
\end{proof}

\begin{example}[\(\Mat_k \simeq \FDVect_k\)]
  We have already met the category \(\Mat_k\) of matrices over a field
  \(k\) (Example~\ref{example:Mat}). Consider the functor
  \[
    F : \Mat_k \to \FDVect_k
  \]
  that maps \(n \in \mathbb{N}\) to the vector space \(k^n\) and matrices
  \(A : m \to n\) (that is, matrices of type \(n \times m\)) to the linear map
  \(x \mapsto Ax\). A basic knowledge of {\sc Linear Algebra} is enough to see
  that \(F\) is full, faithful and essentially surjective on objects. In
  particular, if \(f : k^m \to k^n\) is linear, then the unique matrix
  \(A\) of type \(n \times m\) such that \(FA = f\) is obtained as follows:
  \begin{quote}
    Take the canonical base \(\left\{ e_1, \dots, e_m \right\}\) of
    \(k^m\): then \(A\) is the matrix whose \(i\)-th column
    (\(i \in \left\{ 1, \dots, m \right\}\)) is \(f(e_i)\).
  \end{quote}
  Well, this is the case of the Proposition above: we have a functor
  \(G : \FDVect_k \to \Mat_k\) and natural isomorphisms
  \(\eta : \id_{\Mat_k} \Rightarrow GF\) and
  \(\epsilon : FG \Rightarrow \id_{\FDVect_k}\) such that equations~\eqref{eqn:TriId1}
  and~\eqref{eqn:TriId2} hold. As experiment, let us carry out the
  constructions of the previous proof for this specific case.

  What is \(G\) here? It is any functor such that \(FGV \cong V\): define
  \(GV\) to be \(\dim V\), the dimension of \(V\), for example. What
  about \(\epsilon_V\)? You can choose any linear isomorphism
  \(k^{\dim V} \to V\). What about the action of \(G\) on linear maps now?
  For \(g : V \to W\), the matrix \(Gg\) is determined by
  \(F(G g) = \inv\epsilon_W g \epsilon_V\): in fact,
  \(\inv\epsilon_W g \epsilon_V : k^{\dim V} \to k^{\dim W}\) and it is quite easy to
  reconstruct the matrix \(Gg\) of type \(\dim W \times \dim V\). On the
  other hand, recall this:
  \begin{quote}
    Let \(V\) and \(W\) be vector spaces together with bases
    \(\mathcal{B}_V := \left\{ v_1, \dots, v_m\right\}\) and
    \(\mathcal{B}_W := \left\{ w_1, \dots, w_n\right\}\) respectively and a linear
    map \(f : V \to W\). Write \(p_{\mathcal{B}_V} : V \to k^m\) and
    \(p_{\mathcal{B}_W} : W \to k^n\) the linear maps that send each vector to their
    coordinates in the given base. Then there is a unique matrix \(A\)
    of type \(n \times m\) such that commutes the following diagram
    \[
      \begin{tikzcd}
        V \ar["{p_{\mathcal{B}_V}}", d, swap] \ar["{f}", r] & W \ar["{p_{\mathcal{B}_W}}", d] \\
        k^m \ar["{x \mapsto Ax}", r, swap] & k^n
      \end{tikzcd}
    \]
  \end{quote}
  Well, there is a one-to-one correspondence between the linear
  isomorphisms \(\epsilon_V : k^{\dim V} \to V\) and the bases of \(V\); in that
  case, \(\epsilon_V\) is just the inverse of \(p_{\mathcal{B}_V}\) for one suitable base
  \(\mathcal{B}_V\) of \(V\).

  In substance, the equivalence \(F : \Mat_k \to \FDVect_k\) embodies the
  mantra of your {\sc Linear Algebra} course:
  \begin{quote}
    Once you assign bases, linear maps of finite-dimension vector spaces
    are just matrices.
  \end{quote}

  But let us return to the constructions. What is \(\eta\)? Its components
  are just \(\eta_n : n \to n\), that is for each \(n \in \mathbb{N}\) a
  \(n \times n\) matrix. Look at the proof more closely. \(\eta_m\) is such that
  \(F(\eta_m) = \inv\epsilon_{Fm} = \inv\epsilon_{k^m}\) and
  \(\epsilon_{k^m} : k^m \to k^m\), so it is a very precise matrix: it is the
  matrix that lets you to pass from the canonical base of \(k^m\) to
  another one.
\end{example}



\section{Exercises}

\begin{exercise}
  Given functors \(F_1 : \mathcal{C}_1 \to \mathcal{D}_1\) and \(F_2 : \mathcal{C}_2 \to \mathcal{D}_2\), construct
  a functor
  \[
    F_1 \times F_2 : F_1 : \mathcal{C}_1 \times \mathcal{C}_2 \to \mathcal{D}_1 \times \mathcal{D}_2
  \]
  defined objects as \(F_1 \times F_2 (x, y) := \left( F_1 x, F_2 y \right)\).
\end{exercise}

\begin{exercise}
  Fix \(C \in \obj{\Set}\) and consider the functions
  \begin{align*}
    & \alpha_{A,B} : \Set(C, A) \times \Set(C, B) \to \Set(C, A \times B) \\
    & \alpha_{A,B}(f, g) :=  \lambda x . (f(x), g(x))
  \end{align*}
  for \(A, B \in \obj{\Set}\).
  \begin{enumerate}
  \item These functions are a bijection. More is true though: the
    exercise for you is to understand that \(\eta_{A, B}\) is natural in
    \(A\) and \(B\). To do so, you will need to understand who are the
    functors \(F, G : \Set \times \Set \to \Set\) defined on objects as
    \(F(A, B) := \Set(C, A) \times \Set(C, B)\) and
    \(G(A, B) := \Set(C, A \times B)\).
  \item Prove that for every \(C\) there is a function
    \(\eta_C : C \to C \times C\) which is natural in \(C\) and enjoys the
    following property
    \begin{quote}
      For every \(h : C \to A \times B\) there are unique
      \(f : C \to A\) and \(g : C \to B\) which make the following diagram
      commute
      \[
        \begin{tikzcd}
          C \ar["{\eta_C}", r] \ar["{h}", dr, swap] & C \times C \ar["{f \times g}",
          d] \\
          & A \times B
        \end{tikzcd}
      \]
    \end{quote}
  \end{enumerate}
\end{exercise}

\begin{exercise}
  {\color{red} [Functors of the adjunctions...]}
\end{exercise}

\begin{exercise}
  Let \(\mathcal{C}\) and \(\mathcal{D}\) be two categories, \(\Phi : \op{\mathcal{C}} \times \mathcal{D} \to \Set\) a
  functor and \(F_o : \obj{\mathcal{C}} \to \obj{\mathcal{D}}\) a function. Assume also that
  for every \(a \in \obj{\mathcal{C}}\) we have a natural isomorphism
  \[
    \Phi(a, -) \cong \mathcal{D}(F_o(a), -) .
  \]
  Show that there is a unique functor \(F : \mathcal{C} \to \mathcal{D}\) such that:
  \begin{enumerate}
  \item \(F(a) = F_o(a)\) for every \(a \in \obj{\mathcal{C}}\).
  \item The bijection \(\Phi(a, b) \cong \mathcal{D}(F(a), b)\) is also natural in \(a\).
  \end{enumerate}
  In particular, find a functor \(\Phi^\sharp : \op{\mathcal{C}} \to [\mathcal{D}, \Set]\) such that
  \(\Phi^\sharp (a) = \mathcal{D}(F(a), -)\) for every \(a \in \obj{\mathcal{C}}\).
\end{exercise}

\begin{exercise}
  Consider the functor
  \(\hom_{\mathcal{C}} : \op{\mathcal{C}} \times \mathcal{C} \to \Set\). We have seen how we can associate to
  each \(a \in \obj{\mathcal{C}}\) the functor
  \(\mathcal{C}(a, -) : \mathcal{C} \to \Set\). Well, the exercise for you is: find a functor
  \[
    \yo_{\mathcal{C}} : \op{\mathcal{C}} \to [\mathcal{C}, \Set]
  \]
  such that on objects is defined by
  \(\yo_{\mathcal{C}} (a) := \mathcal{C}(a, -)\). Can you find a functor
  \(\mathcal{C} \to [\op{\mathcal{C}}, \Set]\)? Are the these two faithful? Are they full?
\end{exercise}

\begin{remark}
  The functors above are the {\em Yoneda embeddings}. While it is
  relatively simple to show they are faithful, fullness requires a non
  trivial work to be done, which will be the subject for one of the next
  chapters.
\end{remark}

\begin{figure}
  \centering
  \input{./tikz/el_cat.qtikz}
  \caption{Exercise~\ref{exercise:RepresentableFunctorsInitialObjects}}
\end{figure}

\begin{exercise}\label{exercise:RepresentableFunctorsInitialObjects}
  \begin{enumerate}
  \item Let \(\mathcal{C}\) be a {\em locally small} category and
    \(X : \mathcal{C} \to \Set\) a functor. Assume there is
    \(a \in \obj{\mathcal{C}}\) such that there is a natural isomorphism
    \[
      \psi : \mathcal{C}(a, -) \Rightarrow X
    \]
    Let us introduce \(\tld a := \psi_a\left(\id_a\right)\), which is an
    element of \(X(a)\). Prove that
    \begin{quote}
      for every \(x \in \obj{\mathcal{C}}\) and \(c \in X(x)\) there is a unique
      \(f : a \to x\) such that \(X(f)\left(\tld a\right) = c\).
    \end{quote}
    
  \item Consider a natural isomorphism
    \(\psi : \mathcal{C}(a, -) \Rightarrow \mathcal{C}(b, -)\). Look at
    \(\psi_a(\id_a) : b \to a\) and observe that there is a unique
    \(f : a \to b\) such that \(f \psi_a(\id_a) = \id_b\) (why?). Show that
    also \(\psi_a(\id_a) f = \id_a\). In particular, you can conclude that
    if \(\mathcal{C}(a, -) \cong \mathcal{C}(b, -)\), then \(a \cong b\).
  \end{enumerate}
\end{exercise}

\begin{remark}
  If you have made to this point, congratulations! You have proved one
  of the most important theorems of {\sc Category Theory}, without
  passing before the {\em Yoneda Lemma} -- which is pretty the
  standard. (We will follow this path in one of the next chapters and
  meet some immediate consequences, like the one in the exercise above.)
\end{remark}



%%% Local Variables:
%%% mode: LaTeX
%%% TeX-master: "../CT"
%%% TeX-engine: luatex
%%% End:
